\pagebreak

\section{Simple UART Library File - uarts.3pl}

    \index{library!uarts.3pl}
    This library provides a simple Universal Asynchronous Receiver/Transmitter
    whose configuration is set at compile time. For a UART which can be
    dynamically reconfigured see the previous chapter "General UART library
    file - uart.3pl". The number of data bits is determined from the size
    of the data input or output argument.
    
    This library must be incorporated by using a {\bf module} statement -
    \begin{lstlisting}[gobble=4, language=threepl]
       module "uarts.3pl"
    \end{lstlisting}
    
    For a baud rate $br$ and a current clock frequency $f_c$,
    $br < f_c / 32$ unless $br = f_c / 16$.
    Serial data sampling is controlled by a counter which runs at approximately
    16 times the baud rate. The error on the rate of this counter results
    in a cumulative error which has most impact on the sampling of the
    last serial data bit. If the sampling tolerance exceeds 1/16 of the bit
    width of the last bit a warning message is printed. The calculated sampling tolerance
    includes both the error in the sampling counter and an allowance of
    plus or minus 100ppm in the clocks at either end.


    \subsection{library modules}



        \subsubsection{subsecond\_clocks - millisecond and microsecond
                                    clocks}\label{sec:lmssubsecondclocks}

            {\bf Library: uarts.3pl}

            {\bf Prototype:} \\
            \vsp
            \begin{lstlisting}[gobble=12, language=threepl]
               global module
               subsecond_clocks (value({ms, us}, "log") <- ) {}
            \end{lstlisting}

            {\bf Output parameters:} \\
            {\bf ms} is a millisecond clock.\\
            {\bf us} is a microsecond clock.

            {\bf Input parameters:} \\            
            none           

            {\bf description:} \\
            Outputs a millisecond clock and a microsecond clock.




        \subsubsection{uart\_rx - simple UART receiver}\label{sec:lmsuartrx}

            {\bf Library: uarts.3pl}

            {\bf Prototype:} \\
            \vsp
            \begin{lstlisting}[gobble=12, language=threepl]
               global module
               uart_rx (
                   queue(out), value({oferr, operr, overflow, rts}, "log")
               <-
                   value(rxd, "log"),
                   int(br),
                   value(rxen, "log", true),
                   log(parenb),
                   log(parodd, true),
                   log({cstopb, tq})
               ) {}
            \end{lstlisting}

            {\bf Output parameters:} \\
            {\bf out} is the 8-bit output output.\\
            {\bf oferr} indicates a framing error.\\
            {\bf operr} indicates a parity error.\\
            {\bf overflow} indicates an output queue buffer overflow.\\
            {\bf rts} is 'ready to send'.

            {\bf Input parameters:} \\
            {\bf rxd} is the serial input.\\
            {\bf br} is the baud rate.\\
            {\bf rxen} is receive enable.\\
            {\bf parenb} is parity enable.\\
            {\bf parodd} selects odd parity.\\
            {\bf cstopb} expects a second stop bit.

            {\bf description:} \\
            A UART receiver section. {\bf rxd} is the receive input, usually
            connected direct to an input pin.

            {\bf out} is the output data. The size of the argument to
            this parameter determines the number of bits to be received.

            {\bf oferr} is true if there was a framing error, and {\bf operr}
            indicates a parity error (can only occur if {\bf parenb} is true).
            {\bf oferr} and {\bf operr} are valid whenever out is available.

            Optional input {\bf rxen} if false will disable the receiver when
            the current character is complete. Output {\bf rts} is true when
            the receiver will accept a character. It goes false when {\bf rxen}
            is false or a received character was not acknowledged before the
            expected start of the next character.

            {\bf br} is the baud rate.
            
            If {\bf parenb} is true, an even parity bit is added, odd if {\bf
            parodd} is also true. If {\bf cstopb} is true, a second stop bit is
            added.
            
            {\bf tq} if true shifts the data sampling point from the
            middle of the serial bit period to the 3/4 point. This is
            useful where an open collector driver is the serial data source
            as the -ve transition of the start bit and any data bits is fast
            but +ve transitions rely on the pullup resistor and may be slow.
            Thus the start bit starts the data sampling at the right time
            but later slow +ve transitions on data bits dictate later sampling.

            Flow-control ({\bf rts}) is supported but has not been tested.
            BREAK should be detect but is not.





        \subsubsection{uart\_tx - simple UART transmitter}\label{sec:lmsuarttx}

            {\bf Library: uarts.3pl}

            {\bf Prototype:} \\
            \vsp
            \begin{lstlisting}[gobble=12, language=threepl]
               global module
               uart_tx (
                   value({txd, txbusy}, "log"),
                   value(brc, "log")
                   <-
                   value(in),
                   int(br),
                   value(txen, "log", true),
                   value(cts, "log", true),
                   log(parenb),
                   log(parodd, false),
                   log(cstopb)
               ) {}
            \end{lstlisting}

            {\bf Output parameters:} \\
            {\bf txd} is the serial output.\\
            {\bf brc} is the baud rate clock.

            {\bf Input parameters:} \\
            {\bf in} is the input.\\
            {\bf br} is the baud rate.\\
            {\bf txen} is transmit enable.\\
            {\bf cts} is 'clear to send'.\\
            {\bf csize} is the number of data bits.\\
            {\bf parenb} is parity enable.\\
            {\bf parodd} selects odd parity.\\
            {\bf cstopb} adds a second stop bit.

            {\bf description:} \\
            A UART output section. {\bf txd} is the transmit output, usually
            connected direct to an output pin.

            {\bf in} is the input data. The size of {\bf in} determines the
            number of bits transmitted. If the data input is a queue or an
            expression containing a queue, transmission will only occur when the
            input is available.

            Transmission occurs only if {\bf txen} and {\bf cts} (if connected)
            are asserted. If either go false the transmitter will stop at the
            end of the current character. {\bf cts} usually connects to an
            output pin. {\bf txbusy} indicates the transmitter is operating.

            {\bf br} is the baud rate.
            
            If {\bf parenb} is true, an even parity bit is added, odd if {\bf
            parodd} is also true. If {\bf cstopb} is true, a second stop bit is
            added.

            {\bf brc} is the internal baud clock and is brought out in case
            anything wants it.

            The default is 8-bit characters, no parity, 1 stop bit.

            Flow-control ({\bf cts}) has not yet been tested. {\bf uart\_tx}
            cannot produce BREAK.


    \subsection{library procedures}
    
        None.
        


    \subsection{library functions}
    
        None.


    \subsection{library classes}
    
        None.
